\documentclass[10pt,a4paper]{amsart}
\usepackage[margin=19mm]{geometry}
\usepackage{amsmath,amssymb,mathtools}
\usepackage[scr=rsfs,cal=euler]{mathalfa}
\usepackage{xfrac}
\usepackage[unicode,hidelinks,bookmarks=false]{hyperref}
\newcommand{\ie}{i.e.\ }
\newcommand{\cf}{cf.\ }
\newcommand{\resp}{respectively\ }
\newcommand{\Subsection}{\subsection}
\providecommand{\nobreakdash}{\nobreak\hbox{-}}
\setlength{\emergencystretch}{2em}
\multlinegap0pt
\usepackage{bm}
\usepackage{bbm}
\usepackage{dsfont}
\usepackage{xspace}
\usepackage{cancel}
\usepackage{mathtools}
\usepackage{amsthm}
\makeatletter
\@ifundefined{Prop}{\newtheorem{Prop}{Proposition}}{}
\makeatother
\DeclareMathOperator{\supp}{supp}
\newcommand{\F}[1]{\widehat{\mathbf 1}_{#1}}
\title[On number of cyclic $n$-roots and disjointness of Fourier su]{On number of cyclic $n$-roots and disjointness of Fourier supports}
\author{Weiqi Zhou}
\date{Source: arXiv:2608.06335v2 (CC BY 4.0)}
\begin{document}
\maketitle

\noindent\emph{Source and licence.} On number of cyclic $n$-roots and disjointness of Fourier supports. Version arXiv:2608.06335v2 (CC BY 4.0), distributed under the Creative Commons Attribution 4.0 International licence, as recorded in the arXiv licence metadata for this exact version and shown on the abstract page https://arxiv.org/abs/2608.06335v2. Full rights notice: licenses/arxiv-2608.06335-CC-BY-4.0.txt.\par\smallskip

\noindent\textit{Editorial scope.} Counted unit: the Prop whose source label is \texttt{PropPQR} in the article (source0.tex lines 140--142, source label \texttt{PropPQR}), delivered together with the article's own complete original proof of it (source0.tex lines 144--180, one \texttt{proof} environment of 37 lines). No mathematics is written, repaired or re-derived: statement and proof are the article's own text line for line. Required context retained uncounted: EqCRR (lines 65--67); EqCRR2 (lines 74--76); EqTM (lines 125--127); EqPerp (lines 131--133); EqPoisson (lines 135--137). Every definition, display and label of those lines is retained unchanged. Omitted from this package and not redistributed: 1--64, 68--73, 77--124, 128--130, 134--134, 138--139, 143--143, 181--416. Nothing inside a retained excerpt is omitted or altered: every retained body is delivered exactly as the article's lines are written, so no omission receipt of any kind accompanies this package. Citations: the retained excerpts carry no citation command at all, so no bibliography entry of the article is transcribed and no reference is redistributed. Reference adaptation: none was needed and none was made; every reference inside the delivered unit resolves inside it, so no unresolved reference or citation is produced. Printed slips kept literal: no printed word, symbol, hypothesis or number of the counted statement or of its proof was altered. Layout and class: the article's own class is \texttt{article} with packages including geometry, authblk, amssymb, amsmath, amsthm, mathtools, enumitem, mathabx, hyperref; the wrapper is the lane's pinned \texttt{amsart} editorial wrapper at 10pt on A4, which restates the theorem-like environments the retained text needs and adds the article's own macro definitions that the retained text needs (\texttt{\textbackslash F}, \texttt{\textbackslash supp}), with extra wrapper packages (bm, bbm, dsfont, xspace, cancel, mathtools). The delivered numbering is the unit's own. Assets: no retained span contains a figure environment, a graphics-inclusion command, a PDF-page inclusion, a table or a TikZ picture; no image file is copied into this package, the package declares no asset dependency and no image file is redistributed with it. Licence: version arXiv:2608.06335v2 of this article is distributed under the Creative Commons Attribution 4.0 International licence, as recorded in the arXiv licence metadata for this exact version and shown on the abstract page https://arxiv.org/abs/2608.06335v2. A full rights notice is delivered at licenses/arxiv-2608.06335-CC-BY-4.0.txt. Characters: every retained excerpt is pure ASCII, so the delivered engine, XeLaTeX with its default Latin Modern text font, reports no missing character.
\begin{equation} \label{EqCRR}
u_iv_i=\hat u_i\check v_{i}=0
\end{equation}

\begin{equation} \label{EqCRR2}
\supp(u)\cap\supp(v)=\supp(\hat u)\cap\supp(\check v)=\emptyset.
\end{equation}

\begin{equation} \label{EqTM}
FTx=M\hat x
\end{equation}

\begin{equation} \label{EqPerp}
H_d^{\perp}=\{a\in\mathbb Z_n: ab \bmod n\equiv 0 \text{ holds for all } b\in H_d\}=d\mathbb Z_m=\{0,d,\ldots,(m-1)d\}
\end{equation}

\begin{equation} \label{EqPoisson}
\F{H_d}=\frac{d}{\sqrt n}\cdot \mathbf 1_{H_d^{\perp}}. 
\end{equation}

\begin{Prop} \label{PropPQR}
Let $p,q,r>1$ be natural numbers that are mutually co-prime to each other, then there exist $u,v\in\mathbb C^{pqr}\setminus\{0\}$ that satisfy \eqref{EqCRR}.
\end{Prop}

\begin{proof}
Set
\begin{align*}
A&=1+H_p=1+qr\mathbb Z_p=\{1,1+qr,\ldots,1+(p-1)qr\}, \\
B&=1+H_q=1+pr\mathbb Z_q=\{1,1+pr,\ldots,1+(q-1)pr\}, \\
C&=H_{pq}=r\mathbb Z_{pq}=\{0,r,\ldots,(pq-1)r\},
\end{align*}
and
$$u=q\mathbf 1_A-p\mathbf 1_B, \quad v=\mathbf 1_C.$$
We claim that $u,v$ satisfy \eqref{EqCRR}. First we observe that 
$$A\cap B=\{1\},$$
since the subgroup of order $p$ and order $q$ only intersect at $0$. Thus we get 
$$\supp(u)=A\cup B.$$
Now every element of $\supp(v)=C$ is divisible by $r$, while every element of $\supp(u)$ gives $1$ mod $r$, hence we conclude
\begin{equation} \label{EqUVDisj}
\supp(u)\cap\supp(v)=\emptyset.
\end{equation} 
Next we observe that $\mathbf 1_A$ and $\mathbf 1_B$ are translations of $\mathbf 1_{H_p}$ and $\mathbf 1_{H_q}$ respectively, i.e.,
$$\mathbf 1_A=T\mathbf 1_{H_p}, \quad \mathbf 1_B=T\mathbf 1_{H_q}.$$
Therefore by first applying \eqref{EqTM}, then followed by \eqref{EqPoisson} we obtain
$$\hat u=q\F{A}-p\F{B}=qM\F{H_p}-pM\F{H_q}=\frac{pq}{\sqrt{pqr}}\cdot M(\mathbf 1_{H_p^{\perp}}-\mathbf 1_{H_q^{\perp}}).$$
By \eqref{EqPerp}, $H_p^{\perp}$ is the subgroup generated by $p$, while $H_q^{\perp}$ is the subgroup generated by $q$, they intersect at the subgroup generated by $pq$, therefore
\begin{equation} \label{EqSuppU}
\supp(\hat u)=\{k\in\mathbb Z_n: k \text{ is divisible by precisely one of } p,q \text{ but not } pq\}.
\end{equation}
Applying \eqref{EqPoisson} again we get
$$\check v=\hat v=\F{H_{pq}}=\frac{pq}{\sqrt{pqr}}\cdot \mathbf 1_{H_{pq}^{\perp}}.$$
Here $\check v=\hat v$ holds since both $v$ and $\hat v$ are even functions. By \eqref{EqPerp}, $H_{pq}^{\perp}$ is the subgroup generated by $pq$, therefore
\begin{equation} \label{EqSuppV}
\supp(\hat v)=\{k\in\mathbb Z_n: k \text{ is divisible by } pq\}.
\end{equation} 
Comparing \eqref{EqSuppV} with \eqref{EqSuppU} we get
\begin{equation} \label{EqFUVDisj}
\supp(\hat u)\cap\supp(\check v)=\emptyset.
\end{equation} 
Together \eqref{EqUVDisj} and \eqref{EqFUVDisj} indicate that the so constructed $u,v$ satisfy \eqref{EqCRR2}, and thus also \eqref{EqCRR}. 
\end{proof}

\end{document}
